
\section{Building Blocks of FOL}

\subsection{Building Blocks of FOL} \cite[][Chapter 23, pp. 190-8]{Mag25}

Recall from our discussion that TFL falls short of accurately depicting some
aspects of arguments and inferences in ordinary English. We repeat our ealier
summary for convenience.

\begin{enumerate}
\item TFL does not support the types of inferences depicted in Aristotles
  syllogisms.
\item TFL insists on `excluding the middle' -- every TFL is either true
  or false, but some adjectives do not quite work like that.
\item TFL comes with a number of `paradoxes' that make the material conditional
  unsatisfying for symbolizing English sentences: one the one hand, negating the
  whole conditional often gives the wrong conclusions about its antecedent
  sentence; on the other hand, negating the consequent sentence opens up the
  possibility of contradiction when the antecedant sentence is false.
\end{enumerate}

We will need to wait until future logic courses to deal with (2) and (3)--
\emph{non-classical logic} and \emph{modal logic}, respectively--but we can
start trying to formalize a system to capture Aristotle's syllogism results
right away. This system is \emph{first-order logic}, also commonly refered to
as \emph{predicate logic}.

Recall that the problem with treating these examples in TFL is that TFL does not
have a facility for making distinctions between atomic sentences, aside from
their truth. More precisely, we cannot capture the relationship between ``all
men'' and ``Socrates'' that allows us to make the following inference.

\begin{argument}
\item[ ] All men are mortal. 
\item[ ] Socrates is a man.
\item[\therefore] Socrates is mortal.
\end{argument}

Nor can we capture the fact that ``is/are mortal'' and ``is a man'' are
something like properties that we can apply to things like ``all men'' and
``Socrates''. With first-order logic, we get some rudimentary tools to depict
these phenomena formally.

Let's start with some a informal review of the phenomena.

Consider first some simple sentences that differ as little as possible.

\begin{argument}
\item[ ] John is reading.
\item[ ] Mary is reading.
\item[ ] Fred is reading.
\item[ ] Susie is reading.
\end{argument}

We can see a pattern here involving two components.

First, there is whatever the expressions \emph{John}, \emph{Mary}, \emph{Fred}
and \emph{Susie} have in common. We can tentatively call these expressions
\emph{names} and note that they refer \emph{directly} to the objects that they
represent (which can include any kind of ``person, place or thing''). For all
intents and purposes, we can think of them as representing ``the objects
themselves'' as opposed to some mental or linguistic intermediates. Names are a
special case of \emph{singular terms} which pick out individuals for reference:
they are special because they refer to \textbf{exactly one} item without
ambuguity. We will talk more about other singular terms when we discuss
quantifiers and variables. Names will be symbolized in FOL using lower case
latin letters ($a,b,c,\dots$).

Second, there is the rest of the sentence.

\begin{argument}
\item[ ] $[\hspace{1cm}]_x$ is reading.
\end{argument}

We will refer to such ``incomplete sentences'' as \emph{predicates} and to the
``gaps'' (here labelled with an $x$) as \emph{variables}. Note that predicates
can be made by introducing gaps into multiple locations of sentences, depending
on how many are appropriate to the predicate (i.e. verb, adjective or common
noun) in question. Consider this one which could have a gap in any of three
locations.

\begin{argument}
\item[ ] $[\hspace{1cm}]_x$ gave a book to Mary.
\item[ ] John gave $[\hspace{1cm}]_x$ to Mary.
\item[ ] John gave a book to $[\hspace{1cm}]_x$.  
\end{argument}

We will call a verb like \emph{give} a \emph{three-place predicate}, and
generally refer to a predicate with $n$ gaps as an \emph{n-place predicate}.
When $n>1$ will will sometimes refer to these predicates as \emph{n-place
relations}. Predicates and relations will be symbolized in FOL using capital
latin letters ($A,B,C,\dots$).

Notice that one of the singular terms used in the last example is not a name:
\emph{a book}. This is where the quantifiers come in. Supposing that the gaps
in our predicates can themselves be labeled with symbols (\emph{variables},
symbolized by lower-case latin letters $\dots,x,y,z$), then we can tentatively
express \emph{John gave a book to Mary} using \emph{quantifiers} to \emph{bind}
these variables. Let $G$ symbolize \emph{give}, $j$ and $m$ symbolize
\emph{John} and \emph{Mary}, respectively, $B$ symbolize \emph{is a book} and
$x$ stand for a variable (i.e. the gaps) in both $G$ and $B$ that is also
\emph{bound} by the quantifier \eexists.

\begin{argument}
\item[ ] $\exists(B(x)\eand G(j,x,m))$
\end{argument}

Pronounced in ordinary English, this FOL sentence says ``there is an x such that
x is a book and John gives x to Mary'' which is more or less what we understand
the original English sentence to mean.

% TODO: order of arguments vs variablility of English "give"

We will discuss binding and variables in
more detail as we formalize the logic. For now, note that there are two distinct
quantifiers in first-order logic.

\begin{definition}[FOL Quantifiers (Informal)]
  \label{def:fol-quantifiers-informal}
  First-order logic has two distinct quantifiers:
  \begin{enumerate}
  \item \textbf{Existential Quantifier (\eexists):} expresses the English
    paraphrase \emph{for at least one x}$\dots$
  \item \textbf{Universal Quantifer (\eforall):} expresses the English
    paraphrase \emph{for all x}$\dots$
  \end{enumerate}
  Every quantifier must be followed by a variable indicating which of the
  predicate gaps it \emph{binds} in the interpretation of a sentence. Assigning
  variables to quantifier and gaps is the main mechanism for disabiguating
  complex English sentences within the FOL framework (that is not already
  available in TFL).
\end{definition}

Adding these quantifiers to the language makes it possible to model a variety
of the inferences explored in classical logic by ancient and medievelscholars.

\begin{definition}[Square of Opposition in FOL]
  Just these two quantifiers and negation are enough to express all four corners
  of Aristotle's \emph{square of opposition}. \\ \\
\twoitems{
  \begin{argument}
    \item[ ] Every A is B.
    \item[ ] $\eforall x(A(x)\eif B(x))$
    \item[ ] 
    \item[ ] 
    \item[ ] Not every A is B.
    \item[ ] $\enot\eforall x(A(x)\eif B(x))$
  \end{argument}
}{
  \begin{argument}
    \item[ ] No A is B.
    \item[ ] $\enot\eexists x(A(x)\eand B(x))$
    \item[ ] 
    \item[ ] 
    \item[ ] Some A is B.
    \item[ ] $\eexists x(A(x)\eand B(x))$
  \end{argument}  
}
\end{definition}

Similarly, we can represent additional English quantifiers such as \emph{only}.

\begin{argument}
\item[ ] Only As are Bs.
\item[ ] $\enot\eexists x(B(x)\eand\enot A(x))$
\end{argument}

One question that arises immediately when dealing with a quantifier like
\eforall is ``every what?'' To answer this question we make reference to a
\emph{domain}.

\begin{definition}[Domain (Informal)]
  \label{def:fol-domain-informal}
  In first-order logic, a \emph{domain} is the collection of objects that we
  are talking about.
  \begin{enumerate}
  \item in general, singular terms refer to objects in the domain
  \item every name must pick out \textbf{exactly} one object in the domain;
  \item each domain item can be picked out by any number of names (including
    zero).
  \end{enumerate}
\end{definition}

As we will see when we begin working through the semantics and natural deduction
systems for FOL, the choice of domains is a crucial determinant of the behavior
of the logic.

Among the effects of the domain on FOL involves the absence of referential
material. There are two prominent cases: when a singular term is used, but
it does not refer to any objects in the domain; when a predicate is used, but
it does not apply to any objects in the domain. An expression of the first kind
is \emph{Santa Claus}; an expression of the second is \emph{even prime greater
than 3}. Both of these phenomena can result in unexplected results. To
amelieroate this we will prohibit names that do not refer and we will give
predicates that do not apply to any objects (\emph{empty predicates}) a special
treatment.

\begin{definition}[Empty Predicates]
  When \ebeta is an empty predicate, any sentence
  $\eforall x(\ebeta(x)\eif\dots)$ is (vacuously) true.
\end{definition}

Additionally, when we restict the domain to just people, we can simplify
$\eforall x(P(x)\eif M(x))$ (\emph{Every person is mortal.}) to just
$\eforall xM(x)$ (\emph{Everyone is mortal.}) But quite often, our domains
will be too rich to support this, including not just people but also numerous
places and artefacts as well.

When we symbolize sentences using names, the identity of the name symbol is
enough to indicate that two sentences are ``about'' the same object.

\begin{argument}
  \item[ ] $A(c)\eand B(c)$
\end{argument}

This can be true even when the sentences are separated by any number of others
that may or may not contain any examples of the name in question. The name,
quite simply, always refers to exactly the same thing just in virtue of being a
specific name. 

But when a quantifier is used, the specific sentences to which the quantifier
should be applied is a more delicate matter. The following sentences all express
slightly different things.

\twoitems{
  \begin{argument}
  \item[ ] If c is A, then c is B.
  \item[ ] For all x, if c is A, then x is B.
  \item[ ] For all x, if x is A, then x is B.
  \end{argument}
}{
  \begin{argument}
  \item[ ] $A(c)\eif B(c)$
  \item[ ] $\eforall x(A(c)\eif B(x))$
  \item[ ] $\eforall x(A(x)\eif B(x))$
  \end{argument}
}


Not only do we need to carefully track the across predicates, we also need to
keep them straight within a single predicate. As we noted above, some predicates
(i.e. n-place predicates) have more than one gap to fill when used in FOL
sentences. The example we presented for the verb \emph{give} is repeated here.

\begin{argument}
\item[ ] $[\hspace{1cm}]_x$ gave a book to Mary.
\item[ ] John gave $[\hspace{1cm}]_x$ to Mary.
\item[ ] John gave a book to $[\hspace{1cm}]_x$.  
\end{argument}

We can consider the relevant three-place predicate with all its gaps to be
something like:

\begin{argument}
\item[ ] $[\hspace{1cm}]_x$ gave $[\hspace{1cm}]_y$ to $[\hspace{1cm}]_z$.
\end{argument}

symbolized in FOL as:

\begin{argument}
\item[ ] $G(x,y,z)$
\end{argument}

The order of $x,y,z$ is meaningful and important, but arbitrary. We can assign
the labels in the order that they appear in the original sentence, but note
that this conventions must be used with care, because the following English
sentence has the same meaning, but presents its gaps in a different order.

\begin{argument}
\item[ ] $[\hspace{1cm}]_x$ gave $[\hspace{1cm}]_z$ $[\hspace{1cm}]_y$.
\end{argument}

The \emph{role} that the terms play in the sentence is more important than their
English order when deciding how to assign terms to the FOL predicate.

When more than one quantifier is used, the process of interpretation can get
somewhat complicated. There are three major reasons for this:

\begin{enumerate}
\item The order that quantifiers occur is relevant to the interpretation of
  FOL sentences.
\item The specific choice of variables that the quantifiers apply to is relevant
  to the interpretation of FOL sentences.
\item The TFL connectives, especially negation, support radically different
  interpretations, depending on how they are taken to interact with quantifiers
  in FOL sentences.
\end{enumerate}

Let's look at each of these in turn. First, note that changing the order of the
quantifiers can change the meaning of the sentence.

\begin{argument}
\item[ ] $\eforall x\eexists y L(x,y)$ = Everyone loves someone. (maybe each
  someone different)
\item[ ] $\exists y\forall x L(x,y)$ = There is someone that everyone loves.
  (maybe Sara Lee) 
\end{argument}

Second, the application of a quantifier can vary. We will use `(', `)', `[' and
`]' to indicate the so-called \emph{scope} of a quantifier--the portion of the
overall sentence to which the quantifier applies. Here we have four distinct
quantifiers ($\forall x,\exists z,\exists y,\exists z$) even though the variable
$z$ is used to symbolize two of them.

\begin{argument}
\item[ ] $\forall x[\exists z(D(z)\eand O(x,z)) \eif
  \exists y(\exists z(D(z)\eand O(y,z)) \eand F(x,y))]$
\end{argument}

Finally, the interaction of negation with quantifier scope is a critical to the
interpretation of quantification phenomena that go beyond the two basic
quantifiers. Recall, for example, how we symbolize the square of opposition. The
flexible use of negation and scope allows us to express the FOL sentences with
whichever of the quantifiers that we prefer.

\twoitems{
  \begin{argument}
    \item[ ] Every A is B.
    \item[ ] No A is not B.
    \item[ ] $\eforall x(A(x)\eif B(x))$
    \item[ ] $\enot\eexists x(A(x)\eand\enot B(x))$
    \item[ ] 
    \item[ ] 
    \item[ ] Not every A is B.
    \item[ ] Some A is not B.
    \item[ ] $\enot\eforall x(A(x)\eif B(x))$
    \item[ ] $\eexists x(A(x)\eand\enot B(x))$
  \end{argument}
}{
  \begin{argument}
    \item[ ] No A is B.
    \item[ ] Every A is not B.
    \item[ ] $\enot\eexists x(A(x)\eand B(x))$
    \item[ ] $\eforall x(A(x)\eif\enot B(x))$
    \item[ ] 
    \item[ ] 
    \item[ ] Some A is B.
    \item[ ] Not every A is not B.
    \item[ ] $\eexists x(A(x)\eand B(x))$
    \item[ ] $\enot\eforall x(A(x)\eand\enot B(x))$
  \end{argument}  
}

We still have more to say about singular terms. One open issue is the question
of how to keep them distinct. As it stands, we know that a name always refers
to the same object in the domain (and that that object is indeed in the domain)
but we cannot guess whether two given names refer to one individual or two, both
alternatives are permitted. Similarly, we have no way of expressing the
so-called \emph{count quantifiers} such as \emph{at least two}.

To tackle both of these problems, we will introduce a special two-place
predicate ($=$) that we will interpret as \emph{identity}. In other words,
when we write $a=b$, we will take that to mean that $a$ and $b$ both refer to
the same individual in the domain. So we can symbolize the simple identity
sentence \emph{Cicero is Tully} with the FOL sentence $c=t$. We can use the
negation of identity sentences to express concepts like \emph{everyone else}: 
$\eforall x(x\neq c\eif A(c,x))$, where $x\neq y$ means $\enot(x=y)$. Likewise,
when we want to symbolize \emph{only}, \emph{at least two} and
\emph{at most two}, respectively, we can use the FOL sentences:

%\item[ ] $(A(c)\eand\eforall y(y\neq c\eif\enot A(y)))$
\begin{argument}
\item[ ] $A(c)\eand\enot\eexists x(x=c\eand A(x))$
\item[ ] $\eexists x(A(x)\eand \eexists y(A(y)\eand x\neq y))$
\item[ ] $\enot\eexists x\eexists y\eexists z[((A(x)\eand A(y)\eand A(z))\eand((x\neq y\eand x\neq z)\eand y\neq z)]$
\end{argument}

When we want to express concepts like \emph{exactly 2}, we need to treat it
as \emph{at least 2 \textbf{and} at most 2}, symbolizing it as a conjuntion of
the two.

\begin{argument}
\item[ ] $\eexists x\eexists y[(A(x)\eand A(y)\eand x\neq y]\eand\enot\eexists x\eexists y\eexists z[((A(x)\eand A(y)\eand A(z))\eand((x\neq y\eand x\neq z)\eand y\neq z)]$
\end{argument}

Finally, there is a type of expression whose interpretation seems to be
sometimes best treated as a singular term and sometimes as a quantifier.
Consider the following example.

\begin{argument}
\item[ ] The dog is waiting to come in.
\end{argument}

There are two important steps in the interpretation of this sentence that center
on the use of the expression \emph{the dog} here. First, we are using it like a
singular term to express a proposition about a particular dog that we have in
our domain. Second, we are using quantificational elements to clarify which of
the domain objects we have in mind: namely, the one and only dog in the domain.

Famously, Bertrand Russell argued that we should merge these two aspects of such
expressions (callled \emph{definite descriptions}) by treating \emph{the} as a
type of quantifier. We will look at the other appoach first.

Suppose we just stipulate that \emph{the dog} works like a special kind of name,
picking out its referent ``directly'' and operating as a singular term whenever
it occurs in a sentence. Then we might be able to symbolize the sentence about
our waiting dog as:

\begin{argument}
\item[ ] $W(\iota x(D(x)))$
\end{argument}

Here the expression $\iota x(D(x))$ behaves as a singular term, filling the
gap in the predicate $W$. For Russell's approach, we can just treat \emph{the}
as an \emph{existential quantifier} with a \emph{uniqueness condition}. We've
already seen both of these devices. We can use them together to express our
example:

\begin{argument}
\item[ ] $\eexists x((D(x)\eand\enot\eexists y(x\neq y\eand D(x)))\eand W(x))$
\end{argument}

There are problems with both approaches. We will adopt Russell's, since a full
survey of the $\iota$ operator is beyond the scope of this introduction.

%\begin{definition}[Components of First-Order Logic (Informal)]: There are four
%  basic building blocks in first-order logic. As with TFL, we use familiar
%  characters to \emph{symbolize} the atomic components of the language.
%  \begin{enumerate}
%  \item Connectives (\enot, \eand, \eor, \eif, \eiff): all the same
%    connectives as TFL
%  \item Name ($a,b,c,\dots$): a name refers to exactly one specific thing
%    \emph{directly}, so to speak
%  \item Variable ($\dots,x,y,z$): a variable is another kind of singular
%    term that can be used more abstractly in coordination with quantifiers
%  \item Predicate ($A,B,C,\dots$): a predicate represents an incomplete
%    sentence, i.e. a sentence with \emph{gaps} that can be filled by
%    \emph{singular terms} like names and quantifiers
%  \item Quantifier ($\exists,\forall$): a quantifier is used with variables
%    to say things like ``at least one'' or ``every'' 
%  \end{enumerate}
%\end{definition}

